\section{Topological languages, bases, and preservation}
\label{sec:p03-topological-models-annex}

\begin{annexitem}{annex:p03-topological-polarity}{thm:p03-topological-polarity}{Membership polarity induction}
\begin{proof}
Expand open equality before the induction. A positive membership atom
\(t\varepsilon U\) stays true under enlargement; its negation stays
true under shrinking. Point equality and point predicates do not depend
on \(U\). For conjunction use the two inductions in the same direction.
Negation exchanges the directions by contraposition of the induction
for its matrix. Point and open quantifiers preserve the implication by
using the same witnesses for existential quantifiers and the same
quantification domain for universal quantifiers. Rename their bound variables
so none is the free parameter being changed. This proves the assertion
for all formulas whose free occurrences have the tested polarity, including
guarded ones viewed as their expanded ordinary formulas. It does not require
the parameter values themselves to be members of the quantification domain.
\end{proof}
\end{annexitem}

\begin{annexitem}{annex:p03-topological-basis}{thm:p03-topological-basis}{Topology generation and the guarded-quantifier steps}
\begin{proof}
The empty union is empty and the basis cover gives \(X\).
Arbitrary unions of basis unions are basis unions. To intersect two basis
unions, take the union of every basis element contained in an intersection
of one element from each union. Refinement shows that this union contains
every point of their intersection, and the reverse containment holds by
construction. Thus binary intersections, hence finite intersections,
belong to the generated topology.

For invariance use structural induction allowing arbitrary fixed subset
values for free open variables. Membership, point equality, predicates and
truth have identical interpretations; connectives and point quantifiers
have identical semantics and ranges. Suppose an existential guard
\(\exists U(x\varepsilon U\land\alpha)\) is true in the generated
topology. Its witness is an open \(O\) containing the assigned point.
Some basis element \(V\) contains that point and lies within \(O\).
The negative-polarity lemma makes \(\alpha\) true with \(V\) in place
of \(O\), evaluated still in the generated topology. The matrix induction
then makes it true in the basoid model. Hence \(V\) witnesses its guard
there. Conversely a basis witness is an open witness, and the matrix
induction transfers its truth back. For the universal positive guard,
its negation is the existential negative guard with negated matrix, so the
already proved step and classical negation establish both directions.

For a separating \(L_2\)-example take two points \(a,b\), basis
\(\{\varnothing,\{a\},\{b\}\}\), and its discrete generated topology.
The sentence \(\exists U\forall x(x\varepsilon U)\) is false for
that basis and true for the topology. Moreover, with free parameter
\(V=\{a\}\), the formula \(\forall U(a\varepsilon U\to U=V)\)
is true for the basis and false for the topology. Its equality expansion
has both polarities in \(U\); the definition correctly excludes it from
the guarded positive fragment. Thus treating open equality as a primitive
positive atom would invalidate the proved invariance theorem.
\end{proof}
\end{annexitem}

\begin{annexitem}{annex:p03-topological-compactness-ls}{thm:p03-topological-compactness-ls}{Ordinary first-order basis models and generated topologies}
\begin{proof}
First show that a two-sorted structure satisfying the four basis axioms is
isomorphic to a basoid model. To each open-sort element \(u\) attach its
point extent \(\{x:x\varepsilon u\}\). Ext makes this association
injective and reflects open equality. Union covers the point sort, Empty
supplies its empty extent, and Bas is exactly the local intersection
refinement condition. These extents are a basis. The point predicates keep
their original interpretations, so the extent map gives an isomorphism of
two-sorted structures, not just a surjection onto some topology.

For finitely satisfiable \(\Gamma\), every finite subset of
\(\Gamma\) together with the four axioms has a model: use any
topological model satisfying that finite subset and take its full topology
as a basis. Ordinary two-sorted first-order compactness gives a model of
all of \(\Gamma\) and the four axioms. Equivalently one may use the
standard first-order many-sort encoding, with nonempty disjoint sort
predicates, sort-restricted quantifiers and appropriately typed relations;
its models and satisfaction correspond by term/formula induction.
The extent construction gives a basoid model. Generate its topology and
apply basis invariance to each sentence in \(\Gamma\). This produces
the required topological model.

For countable signature apply ordinary downward L\"owenheim--Skolem to
the two-sorted structure associated with a given topological model. A
countable elementary substructure still satisfies all of \(\Gamma\)
and the four axioms. Its point sort and its open sort are countable.
Take extents in its point sort, giving a countable basis on a countable
carrier, then generate the topology. Basis invariance preserves every
sentence of \(\Gamma\). For an infinite-point model add the true
sentences asserting at least \(n\) distinct points, for all finite \(n\),
before the construction; the resulting countable carrier is infinite.
Free point/open parameters in a satisfiable family of formulas can instead
be named by constants in their respective sorts, with their assignments
kept fixed. The countable conclusion requires the expanded signature to
remain countable. Generation adds possibly uncountably many unions; no
bound on the cardinality of the full topology follows.
\end{proof}
\end{annexitem}

\begin{annexitem}{annex:p03-topological-property-formulas}{thm:p03-topological-property-formulas}{Meaning, polarity and the Alexandroff direction}
\begin{proof}
For distinct points the \(T_0\) formula chooses an open containing
one but not the other, which is exactly distinguishability by opens.
The \(T_1\) formula makes this choice in each ordered direction,
equivalently each singleton is closed: if \(\{y\}\) is closed,
its open complement contains every \(x\ne y\); conversely the union
of the chosen neighborhoods of all such \(x\) is that complement.
The \(T_2\) guard chooses a neighborhood for each point and its matrix
says their intersection is empty, exactly the Hausdorff separation axiom.
The discreteness formula gives an open neighborhood contained in \(\{x\}\)
and containing \(x\), hence the singleton itself. Every union of those
singletons is open. Conversely a discrete topology supplies them.

The Alexandroff sentence gives \(U_x\) contained in every neighborhood
of \(x\). If arbitrary intersections of opens are open, their intersection
over all neighborhoods of \(x\) supplies just such a \(U_x\).
Conversely consider any intersection \(A\) of opens. Every \(x\in A\)
has its least neighborhood \(U_x\) contained in every member of that
family, hence in \(A\). Thus \(A=\bigcup_{x\in A}U_x\) is open.
For an empty family the intersection is \(X\), already open. With the
reversed inclusion, \(X\) would always be a permissible witness; the
ordinary real line, in which a point has no least open neighborhood,
separates that erroneous direction from the actual property.
In the displayed formulas each existentially bound neighborhood occurs
only negatively in its matrix, while each universally bound neighborhood
occurs positively. Their guarding membership is outside that matrix.
Expansion of abbreviations therefore meets the \(L_t\) polarity rules.

For connectedness the antecedent says \(V=X\setminus U\), so both
\(U\) and its complement are open. The conclusion says \(U=X\) or
\(U=\varnothing\). This excludes exactly nontrivial clopen subsets.
It proves the asserted semantics in full \(L_2\), without pretending
that unguarded open quantification belongs to \(L_t\).
\end{proof}
\end{annexitem}

\begin{annexitem}{annex:p03-topological-translation-homeomorphism}{thm:p03-topological-translation-homeomorphism}{Two satisfaction inductions}
\begin{proof}
For the modal translation, truth, propositional letters, conjunction and
negation follow directly from their matching semantics and induction.
The box step says the assigned point belongs to the interior of the
matrix truth set. This holds exactly when there is an open neighborhood
of that point contained in the truth set. The matrix induction translates
containment into \(\forall y(y\varepsilon U\to\mathrm{ST}_y(\varphi))\).
This is precisely the translated existential guard. Fresh variable choices
ensure no free parameter is captured.

For homeomorphism invariance let \(h:X\to X'\) be the specified map.
It gives a bijection between points and a bijection between opens by
\(U\mapsto h[U]\). Point equality is preserved and reflected by
injectivity; open equality is preserved and reflected by the induced
set bijection; membership is preserved and reflected because
\(x\in U\Longleftrightarrow h(x)\in h[U]\). Predicate atoms agree
by the assumption on their interpretations. Negation and conjunction
follow by induction. An existential point or open witness is transported
by its corresponding bijection; any target witness has a preimage, proving
the converse. Universal quantification follows either by the same bijections
or by negated existential quantification. This proves all \(L_2\) formulas.
An ordinary structure isomorphism without the topological requirement
need not transport neighborhoods. For example the identity on a two-point
pure equality carrier is an ordinary isomorphism from its indiscrete to
its discrete topology, but only the latter satisfies the discreteness
sentence. This gives the stated failure boundary explicitly.
\end{proof}
\end{annexitem}

\begin{annexitem}{annex:p03-topological-rip-structure}{thm:p03-topological-rip-structure}{The quotient topology and canonical box factors}
\begin{proof}
Full products are nonempty by Choice. Element equivalence is reflexive
and symmetric; its transitivity follows by intersecting the two filter-large
equality sets. For finitely many changed input representatives, their common
equality set is filter-large. On it functions give identical outputs and
relations have identical truth. Upward closure of \(\Psi\) therefore
makes all interpretations independent of representatives. Term induction
gives coordinatewise term classes.

Whole product belongs to \(B_\Lambda\), and the intersection of two
boxes has its support contained in the union of their two supports, hence
again in \(\Lambda\). This proves the ordinary basis property.
For a nonempty box \(B=\prod_i U_i\), fix \(b\in B\).
Its saturation under element equivalence is
\[
 q^{-1}[q(B)]=\bigcup_{F\in\Psi}\prod_i W_i^F,
 \qquad W_i^F=\begin{cases}U_i&i\in F,\\S_i&i\notin F.\end{cases}
\]
One containment follows from agreeing with a member of \(B\) on a
large set. Conversely for a tuple in the \(F\)-box, use its coordinates
on \(F\) and those of \(b\) elsewhere; this tuple belongs to \(B\)
and is equivalent to it. Each displayed box has support contained in the
original support. The saturation is open. For an empty box the saturation
is empty, also open. Images of arbitrary open unions of boxes are therefore
open in the quotient topology, so \(q\) is open. If a quotient open
contains \([f]\), its preimage contains \(f\) in some basic box;
the image of that box is a neighborhood of \([f]\) contained in the
quotient open. These images are thus a basis.

For a nonempty box Choice, or the already fixed \(b\), lets any member
of a factor \(U_i\) extend to a member of the box. Its canonical
projection is therefore the original factor set. The saturation computation
proves the membership criterion with these factors. For an empty box all
projections are empty, so its local membership truth set is empty and
does not belong to the proper filter; both sides of the criterion are false.
This explains why noncanonical factors of an empty written rectangle cannot
be used for local parameters. On \(I=\mathbb N\) with the cofinite
filter, the rectangle empty only at coordinate zero is empty, while its
noncanonical coordinate-membership truth set would be cofinite. Using its
actual projections avoids that false atomic conclusion.

For the Boolean quotient, equivalence of sets is agreement on a
filter-large set; finite intersections of those agreement sets prove
transitivity and compatibility with union, intersection and complement.
Thus the source's declared operations are well defined and satisfy the
Boolean identities. Properness makes \([\varnothing]\ne[I]\).
The image \(\Lambda/{\equiv}\) contains zero and is closed under
finite joins. If \([A]\le[L]\) with \(L\in\Lambda\), then
\([A]=[A\cap L]\), and \(A\cap L\in\Lambda\); this proves
downward closure. So \(\mathbf B\) satisfies \(T_B\).

With the filter \(\{I\}\), equivalence is actual identity and the
quotient map is the identity. Finite supports give the usual product
topology; arbitrary supports give the box topology. For infinitely many
two-element discrete factors, the full box topology is discrete since each
point has its singleton box. The product topology is not: any basic
neighborhood restricts only finitely many coordinates and a different
tuple can be obtained by changing another coordinate. Thus the topology
distinction has an explicit separating example.
\end{proof}
\end{annexitem}

\begin{annexitem}{annex:p03-topological-carrier-substitution}{thm:p03-topological-carrier-substitution}{Free occurrences and the whole-open value}
\begin{proof}
For an atom involving the designated free variable, assigning \(Y=S\)
makes \(t\varepsilon Y\) true exactly as \(t=t\) is true. Other
atoms do not change. Negation and conjunction preserve the equivalence
by induction. Quantifiers over point variables or over other open variables
use the same ranges on both sides; their assignments commute with changing
\(Y\), after fresh bound names are chosen. A quantifier binding \(Y\)
has no free designated occurrences in its scope, which therefore remains
unchanged. These cases prove the full structural induction. Replaced atoms
introduce no open occurrences, so all other free-open polarities remain
the same. The resulting formula has no free \(Y\).
\end{proof}
\end{annexitem}

\begin{annexitem}{annex:p03-topological-fv}{thm:p03-topological-fv}{All Boolean clauses and both actual existential witnesses}
\begin{proof}
Induct simultaneously on reduction correctness, increasingness of its
Boolean condition, its free-variable bound and retention of free-open
polarities. Terms give coordinate classes by the quotient-structure proof.
Point atoms therefore use the filter-large truth-set condition, equivalently
\([I_\alpha]=1\). Membership atoms use the proved canonical-box
criterion. Open equality has already been expanded extensionally into point
quantification and membership, and is not an additional atomic case.
Truth and falsity use their declared constant conditions.
These conditions are increasing. Conjunction preserves correctness and
increasingness by its two induction hypotheses and concatenated lists.

For negation, the local truth class of \(\neg\theta_j\) is
\([I_{\theta_j}]'\). Thus evaluation of
\(\neg\tau(y_1',\ldots,y_m')\) at those classes gives exactly
the negation of the old condition. If \(\bar y\le\bar t\), then
\(\bar t'\le\bar y'\). Increasingness of \(\tau\) gives
\(\tau(\bar t')\to\tau(\bar y')\), so contraposition proves
increasingness of the new condition. Negating every local formula also
reverses exactly the input's free-open polarities. Conjunction and all
quantifier clauses do not introduce new free variables, and their point
binders remove just the quantified point slots.

For an individual existential, an actual quotient witness \([g]\)
defines
\(Z_k=\{i:(S_i,O_i)\models\bigwedge_{j\in s_k}\theta_j[f_i,U_i,g_i]\}\).
These sets satisfy \(Z_i\cap Z_j=Z_k\) whenever
\(s_i\cup s_j=s_k\), and each is contained in \(I_{\psi_k}\).
Their singleton classes satisfy \(\tau\) by matrix induction, proving
the forward clause.
Conversely choose representatives \(Z_k\) of Boolean witnesses.
Each of the finitely many inequalities and intersection identities holds
on some \(\Psi\)-large set; intersect these sets into \(F\in\Psi\).
For \(i\in F\) let \(s(i)=\{j\le m:i\in Z_j\}\).
If this set is nonempty and equals \(s_h\), repeated application of
the intersection identities gives
\(i\in Z_h\subseteq I_{\psi_h}\) on \(F\).
There is therefore a factor point satisfying every requested
\(\theta_j\). Choice selects such points for these indices and
arbitrary points for the indices with no requests or outside \(F\).
Every requested singleton truth set contains \(Z_j\cap F\);
its Boolean class is at least \([Z_j]\). Increasingness of \(\tau\)
and matrix induction make \([g]\) a witness. No exact negative-pattern
realization is needed or asserted: this is why the construction keeps
the input's open polarities. The subset conjunctions and their existential
point quantifiers retain those polarities.

Now consider the negative open guard. Matrix induction provides
\((\tau;\theta_1,\ldots,\theta_m)\), with \(Y\) negative in
every \(\theta_j\). Hence all \(\psi_k\) are admitted guards;
whole-carrier substitution gives the \(\eta_j\), retaining every other
open role. The new condition is increasing: increasing \(y_k\) leaves
its upper bounds weaker, and increasing \(v_j\) makes
\(z_j\setminus v_j\) smaller, so the ideal axiom preserves
\(\lambda(z_j\setminus v_j)\). The other conditions do not change.

Suppose first that a neighborhood witnesses the guard in the quotient
topology. Shrink it to an image \(q(V)\) of a basic box containing the
assigned term value; negativity of the matrix preserves its truth.
Basis invariance allows matrix induction at this basic neighborhood.
Since it contains a point, \(V\) is nonempty, so its canonical factors
are genuine opens \(V_i\), equal to \(S_i\) outside some
\(L\in\Lambda\). Define \(Z_k\) as the truth sets of the subset
conjunctions at \((f_i,U_i,V_i)\). Guard membership gives
\(F=\{i:t[f_i]\in V_i\}\in\Psi\); hence
\(Z_k\cap F\subseteq I_{\psi_k}\). Their intersections satisfy
the declared identities, and singleton classes satisfy \(\tau\).
Outside \(L\), whole-carrier substitution makes
\(\theta_j\) and \(\eta_j\) agree. Thus
\(Z_j\setminus I_{\eta_j}\subseteq L\), establishing every
ideal condition and proving the forward implication.

Conversely take Boolean witnesses and representatives \(Z_k\).
The availability inequalities and finite intersection identities hold
exactly on a common large set. For each ideal condition, choose
\(C_j\in\Lambda\) representing
\([Z_j\setminus I_{\eta_j}]\). Equality of these classes holds
on a large set \(G_j\), so
\((Z_j\setminus I_{\eta_j})\cap G_j=C_j\cap G_j\in\Lambda\).
Include all the \(G_j\) in the finite intersection, giving \(F\in\Psi\)
with all properties at once. Put
\[
                    L=\bigcup_{j=1}^m
                       ((Z_j\setminus I_{\eta_j})\cap F)\in\Lambda.
\]
For \(i\in L\), the requested set \(s(i)\) is nonempty.
Its intersection identity and availability bound yield
\(i\in I_{\psi_h}\) for the corresponding \(s_h\).
Choose an actual \(V_i\in O_i\) containing \(t[f_i]\) and
satisfying all those requested \(\theta_j\). For \(i\notin L\)
set \(V_i=S_i\). Choice performs the possibly infinite selection
inside \(L\), and the resulting box \(V\) has support contained
in \(L\), so belongs to \(B_\Lambda\). It contains the term
tuple at every coordinate, hence its quotient image contains the term class.

If \(i\in Z_j\cap F\), then either \(i\in L\) and the chosen
factor neighborhood satisfies \(\theta_j\), or \(i\notin L\)
and \(i\in I_{\eta_j}\); whole-carrier substitution then makes
\(\theta_j\) true at \(V_i=S_i\). Consequently each actual
singleton truth class is above \([Z_j]\). Increasingness of \(\tau\)
and matrix induction prove the matrix at \(q(V)\), giving the
actual guarded witness. This completes both existential cases.
Every list and conjunction is finite; recursion on the finite formula
tree terminates. Universal clauses and remaining connectives are their
declared classical expansions. All asserted correctness, polarity and
monotonicity properties are therefore proved together.
\end{proof}
\end{annexitem}

\begin{annexitem}{annex:p03-topological-transfer}{thm:p03-topological-transfer}{Sentence transfer and the ultrafilter guard induction}
\begin{proof}
For an input sentence its reduction components are sentences: the
free-variable bound proves this at every recursive step. Factorwise
\(L_t\)-equivalence makes their truth sets identical in the two families.
The same coefficient structure and same Boolean condition therefore
give the same truth in the two reduced ideal-products.

For the ultrafilter assertion one can prove the stronger one-component
criterion directly. Atomic point formulas and membership have already
been verified. For negation and conjunction, the ultrafilter alternative
and finite intersection laws give exactly the local truth-set identities.
For an individual existential, a quotient witness makes its local matrix
truth large by induction and hence its local existence truth large.
Conversely choose actual factor witnesses on that large existence set and
arbitrary points elsewhere; their quotient class witnesses the matrix by
induction.

For a negative open guard, a true quotient guard can be shrunk to a
basic open witness using negativity and basis invariance. Its term-membership
and matrix truth sets are both large by induction, so their intersection
makes the factor guard true. Conversely if its local guard is true on a
large set, choose in each such factor an actual open witness containing
the term point and satisfying the matrix; choose the whole factor carrier
elsewhere. With \(\Lambda=\mathcal P(I)\) this is an allowed full box.
Its canonical factors are the chosen nonempty opens. Term membership
holds everywhere, matrix truth on the large set, and induction makes its
quotient image a witness. Dual guarded universal quantification follows
by negation. This proves all formulas, not just sentences.

For a diagonal ultrapower, constant point tuples and constant basic-open
tuples give local truth sets either all indices or empty. The proved
criterion yields elementary preservation and reflection on those
parameters. Unions of quotient basic opens generate its topology; basis
invariance, not an identification of the basis with all opens, justifies
the topological interpretation. For a basic point predicate its extension
is therefore precisely the classes of tuples whose local membership sets
are large. If one factor predicate extension is empty and all others are
full, its Cartesian product is empty but its local truth set may be large
for a nonprincipal ultrafilter. This explicit case proves why using the
empty Cartesian product's image as its interpretation would give an
incorrect atomic clause.
\end{proof}
\end{annexitem}
